\section{Introduction}
  \label{sec:intro}

  The Admissibility Sub-Programme of \textsc{Cosmochrony} asks which spectra of the
  relational operator $\mathcal{L}_\chi$ are compatible with the bounded-flux evolution
  constraint $\|c_{\mathrm{BI}}\| \le 1$.
  The programme proceeds in three steps.
  \begin{itemize}
    \item \textbf{Step~1} (\emph{SpectralAdmissibility~1.0}):
    generators satisfying the neutrality condition $\chi_\rho(s) = 0$ correspond
    to traceless elements of $\SU(2)$, i.e.\ to half-turns.
    \item \textbf{Step~2} (\emph{SpectralCapacity~1.0}):
    the Klein classification of finite subgroups of $\SU(2)$ restricts the
    possible groups $\rho(G)$ to
    $\{\mathrm{Dic}_n\,(n \ge 2),\; 2T,\; 2O,\; 2I\}$.
    \item \textbf{Step~3} (\emph{this paper}):
    pairwise character geometry distinguishes the cases analytically and
    establishes which groups actually arise.
  \end{itemize}

  The key difficulty left open after Step~2 is that the Klein classification is
  typically proved by geometric or group-theoretic means, neither of which refers
  to the spectral constraint.
  The main contribution of the present paper is to show that the full classification
  \emph{follows} from two character-theoretic invariants:
  the second-moment tensor $M$ and the Gram matrix $G_{st}$.
  No geometric input beyond neutrality is needed.

  \paragraph{Main result (informal).}
    The finite subgroups of $\SU(2)$ admissible under spectral neutrality are exactly
    \[
      Q_8,\quad \mathrm{Dic}_n\;(n \ge 3),\quad 2O,\quad 2I.
    \]
    The binary tetrahedral group $2T$ is \emph{structurally excluded}: the full set of
    its traceless elements generates only $Q_8$, not $2T$.
    The three remaining cases are separated by whether $M \propto I$ (isotropy)
    and by whether the Gram matrix contains values outside $\{0, \pm 1\}$.

  \paragraph{Organisation.}
    Section~\ref{sec:setup} recalls the setup and the results of Steps~1--2.
    Section~\ref{sec:gram} establishes the Gram identity $G_{st} = -\tfrac{1}{2}\chi_\rho(st)$.
    Section~\ref{sec:M} analyses the second-moment tensor and proves that
    $\mathrm{Dic}_n$ is anisotropic for $n \ge 3$.
    Section~\ref{sec:Q8} proves that an orthogonal Gram pattern forces $\rho(G) \cong Q_8$.
    Section~\ref{sec:2T} proves the structural exclusion of $2T$.
    Section~\ref{sec:theorem} assembles the main classification theorem.
    Section~\ref{sec:klein} discusses completeness via the Klein list.
    Section~\ref{sec:discussion} draws the conceptual conclusions.
